\section{Introduction}
\label{sec:introduction}

The development of policy optimization has been marked by two representative milestones. The first is the Zero series of DeepMind around 2017, where AlphaGo Zero, AlphaZero, and MuZero~\citep{ref_alphagozero,ref_alphazero,ref_muzero} use Monte Carlo tree search (MCTS) with policy priors and posterior return/value information to produce improved action-visit distributions, then project them back into the policy network, pushing AI systems beyond human-level performance in board games and game-playing domains. We refer to this route as \emph{Search and Policy Projection}. The second is the recent RLVR paradigm exemplified by DeepSeek-R1~\citep{ref_r1}, where large-batch rollouts are used to collect rewards and policy gradients directly update the model, eliciting reasoning abilities in large language models. We refer to this route as \emph{Rollout and Policy Gradient}.

The success of search raises a natural question: why do modern policy-gradient methods still rely primarily on plain rollouts, rather than systematically exploiting posterior information as in AlphaZero-style systems? There are three main obstacles: search changes the sample measure of tree trajectories relative to chain rollouts, inducing distributional shift; classical MCTS is sequential and CPU-intensive, making it difficult to match the GPU-parallel rollout infrastructure of PPO/RLVR; and Policy Projection is aligned with breadth-wise action-visit statistics, whereas Policy Gradient seeks high-return trajectories and thus calls for depth-oriented search that rolls out from intermediate states.

Recent work such as TreePO, TreeRL, SEEA-R1, and ARPO~\citep{ref_treepo,ref_treerl,ref_seear1,ref_arpo} has explored using tree-structured samples for policy-gradient training. However, three key questions remain open. \textbf{(i) How should policy gradients on tree trajectories be aggregated?} Under naive equal-weight aggregation, even if every continuation is sampled from the current policy, repeatedly selecting an existing intermediate state for further sampling makes the suffix trajectories after the branching point contribute a larger relative share to the policy-gradient estimator than they would under a purely on-policy chain rollout. \textbf{(ii) How should advantages be estimated on tree structures?} Standard GAE~\citep{ref_gae} assumes a single chain successor, whereas a tree trajectory contains multiple continuations from the same node, requiring an expectation-consistent multi-successor advantage recursion. \textbf{(iii) How should branching positions be selected?} Existing methods often rely on heuristic signals to decide where to continue expansion, lacking a principled criterion for why search should resume there.

To address these questions, we derive On-policy Parallel Tree Search (OPTS) and Tree Trajectory Policy Optimization (TTPO), systematizing the Search and Policy Gradient line of policy optimization. Our contributions are twofold.\footnote{To support anonymous review and reproducibility, we provide an anonymized code repository in the supplementary material. Code and experimental logs will be released upon acceptance.}

\paragraph{Method.} We propose OPTS-TTPO, a unified Search and Policy Gradient framework in which OPTS performs search and TTPO performs learning. We prove that the branch-weight-corrected policy-gradient estimator is unbiased, introduce Tree-based Generalized Advantage Estimation (TreeGAE), and derive the On-policy Trajectory Rebranching Criterion (OTRC).
\paragraph{Empirical results.} We provide unified implementations and cross-domain evaluations on LLM mathematical reasoning, including both training-time and test-time search, Atari-57 discrete control~\citep{ref_atari}, and MuJoCo continuous control~\citep{ref_mujoco}. OPTS-TTPO outperforms DAPO~\citep{ref_dapo} and PPO~\citep{ref_ppo} on verifiable mathematical reasoning tasks, OPTS improves over majority voting under the same test-time search budget, and OPTS-TTPO achieves stronger overall performance than PPO on Atari-57 and MuJoCo. Variance-reduction experiments further show that, across batch-size settings, positive-advantage samples have lower policy-gradient estimation variance than negative-advantage samples, and OPTS-TTPO consistently yields lower variance than PPO.
